\section{A fixed denominator in the Gamma weight}
\label{sec:lc-fixed-denominator}

Throughout this section the original real number $c=k/2$, the original
coordinate $S$, and the actual monic polynomial
\[
 \chi(S)=\prod_{\lambda\in\Lambda}(S-\lambda)^{m_\lambda},
 \qquad q=\sum_{\lambda\in\Lambda}m_\lambda\ge1,
\]
are fixed. Distinct elements of $\Lambda$ denote distinct roots, and every
$m_\lambda$ is the full multiplicity in $\chi$. No root is discarded. Set
\[
 \sigma(u)=\frac{|\Gamma(1/4+iu/2)|^2}{2\pi},
 \qquad w_M(u)=\frac{\sigma(u)}{(1+u^2)^M},
 \qquad M\in\mathbb Z_{\ge0}.
 \tag{LC.1}\label{eq:lc-weights}
\]
For a positive weight $w$ with finite moments, let
\[
 \|p\|_w^2=\int_{\mathbb R}|p(c+iu)|^2w(u)\,du,
 \qquad p\in\mathcal P_N=\mathbb C[S]_{\le N}.
 \tag{LC.2}\label{eq:lc-norm}
\]
For $N\ge q-1$ and a remainder $a(S)=\sum_{j=0}^{q-1}a_jS^j$, put
\[
 \|[a]\|_{w,N,\chi}^2
 =\min_{\substack{p\in\mathcal P_N\\p-a\in(\chi)}}\|p\|_w^2
 =\boldsymbol a^*G_N(w,\chi)\boldsymbol a,
 \qquad V_N(w,\chi)=\det G_N(w,\chi).
 \tag{LC.3}\label{eq:lc-quotient}
\]
The minimum exists uniquely because it is the minimum of a positive definite
quadratic form on a nonempty affine subspace of a finite dimensional space.
The form is positive definite since a polynomial vanishing almost everywhere
on $c+i\mathbb R$ vanishes identically. In particular, $V_N$ retains the
original Gram determinant convention.

\subsection{Explicit comparison of the weights}

The exact Gamma envelope, including $u=0$, is
\[
 \begin{gathered}
 c_\Gamma e^{-\alpha|u|}(1+|u|)^{-1/2}
 \le\sigma(u)\le
 C_\Gamma e^{-\alpha|u|}(1+|u|)^{-1/2},\\
 \alpha=\frac\pi2,
 \quad c_\Gamma=\sqrt2e^{-7/6},
 \qquad C_\Gamma=\sqrt{2\sqrt5}e^{2/3}.
 \end{gathered}
 \tag{LC.4}\label{eq:lc-gamma-envelope}
\]
For completeness, its proof is recalled. Binet's exact identity on $\Re z>0$,
with the logarithm real on positive reals, is
\[
 \log\Gamma(z)=(z-\tfrac12)\log z-z+\tfrac12\log(2\pi)+R(z),
 \quad
 R(z)=\int_0^\infty e^{-zt}\frac1t
 \left(\frac1{e^t-1}-\frac1t+\frac12\right)dt.
 \tag{LC.5}\label{eq:lc-binet}
\]
This identity is given in NIST DLMF, equation 5.9.10$_2$,
\url{https://dlmf.nist.gov/5.9.E10_2}.
Writing $x=t/2$ in the kernel, its value is
$(\coth x-1/x)/(4x)$. The functions
$F(x)=x\cosh x-\sinh x$ and
$H(x)=(x^2+3)\sinh x-3x\cosh x$ vanish at zero and satisfy
$F'(x)=x\sinh x\ge0$ and $H'(x)=xF(x)\ge0$ for $x\ge0$.
Dividing by the positive denominators $x\sinh x$ and $3x\sinh x$ gives
$0\le\coth x-1/x\le x/3$ for $x>0$. Thus the kernel in
\eqref{eq:lc-binet} lies between $0$ and $1/12$, and
$|R(z)|\le1/(12\Re z)$.

At the retained argument $z=1/4+iu/2$, put
$x=|u|$ and $\rho=|z|=(1/16+x^2/4)^{1/2}$. Twice the real part of
\eqref{eq:lc-binet} gives exactly
\[
 \sigma(u)=\rho^{-1/2}
 \exp\bigl(-x\arctan(2x)-\tfrac12+2\Re R(z)\bigr).
 \tag{LC.6}\label{eq:lc-exact-gamma}
\]
The remainder bound is $|R(z)|\le1/3$. For $x\ge0$ the exponent lies
between $-\alpha x-7/6$ and $-\alpha x+2/3$: the lower inequality uses
$\arctan(2x)\le\pi/2$, while for $x>0$ the upper inequality follows from
$\arctan(2x)=\pi/2-\arctan(1/(2x))$ and
$\arctan y\le y$ for $y\ge0$. At $x=0$ the upper inequality follows
directly from $-1/2\le0$. The exact squared differences
\[
 \frac{(1+x)^2}{4}-\rho^2=\frac{3+8x}{16},
 \qquad
 \rho^2-\frac{(1+x)^2}{20}=\frac{(4x-1)^2}{80}
\]
give $(1+x)/(2\sqrt5)\le\rho\le(1+x)/2$. Combining these inequalities
with \eqref{eq:lc-exact-gamma} proves \eqref{eq:lc-gamma-envelope}, including
the origin.

Fix $r>1$ and put $\beta=\alpha(r-1)>0$. Define the finite explicit constant
\[
 B_{r,M}=\frac{C_\Gamma}{c_\Gamma}
 \sum_{j=0}^{2M}\binom{2M}{j}\frac{j!}{\beta^j}.
 \tag{LC.7}\label{eq:lc-constant}
\]
For $x=|u|\ge0$, the Gamma envelope gives
\[
 \frac{\sigma(ru)}{\sigma(u)}
 \le\frac{C_\Gamma}{c_\Gamma}e^{-\beta x}
 \left(\frac{1+x}{1+rx}\right)^{1/2}
 \le\frac{C_\Gamma}{c_\Gamma}e^{-\beta x}.
\]
For each integer $j\ge0$, the exponential series implies
$e^{\beta x}\ge(\beta x)^j/j!$, so that
$x^je^{-\beta x}\le j!\beta^{-j}$, with the case $j=0$ read directly.
Since $(1+x^2)^M\le(1+x)^{2M}$, expansion of the latter polynomial now proves
\[
 B_{r,M}^{-1}\sigma(ru)\le w_M(u)\le\sigma(u)
 \qquad(u\in\mathbb R).
 \tag{LC.8}\label{eq:lc-pointwise-sandwich}
\]
Every comparison is global, and $r$ is fixed independently of $N$.

\subsection{The exact morphism on the original jet constraints}

Order the roots once and use that ordering throughout. For a polynomial $p$,
write its divided jets as
\[
 J_\chi p=\left(\frac{p^{(j)}(\lambda)}{j!}\right)_
 {\lambda\in\Lambda,\ 0\le j<m_\lambda}.
\]
Let $E_\chi$ be the matrix of $J_\chi$ on the original remainder frame
$1,S,\ldots,S^{q-1}$. This matrix is invertible: a remainder in its kernel
has a zero of order at least $m_\lambda$ at every $\lambda$, is consequently
divisible by $\chi$, and has degree less than $q$, hence is zero.
For all polynomials $p,a$, the same multiplicity criterion proves
\[
 p-a\in(\chi)\quad\Longleftrightarrow\quad J_\chi p=J_\chi a.
 \tag{LC.9}\label{eq:lc-exact-constraints}
\]

For the comparison parameter $r>0$, define the polynomial bijection
\[
 (T_rp)(S)=p\left(c+\frac{S-c}{r}\right),
 \qquad
 \zeta_\lambda(r)=c+r(\lambda-c),
 \qquad
 D_r=\operatorname{diag}_{\lambda,j}(r^j).
 \tag{LC.10}\label{eq:lc-pullback}
\]
Its inverse is $T_{1/r}$. The locations $\zeta_\lambda(r)$ are pairwise
distinct, because $r>0$. Write $J_r$ for the divided-jet map of the same
multiplicities at these locations. The chain rule, with all constants
retained, gives
\[
 \frac{(T_rp)^{(j)}(\zeta_\lambda(r))}{j!}
 =r^{-j}\frac{p^{(j)}(\lambda)}{j!},
 \qquad J_rT_r=D_r^{-1}J_\chi.
 \tag{LC.11}\label{eq:lc-jet-morphism}
\]
The literal substitution $v=ru$ in the integral on the original line gives
\[
 \|p\|_{\sigma(r\cdot)}^2
 =r^{-1}\int_{\mathbb R}
 \left|p\left(c+\frac{iv}{r}\right)\right|^2\sigma(v)\,dv
 =r^{-1}\|T_rp\|_\sigma^2.
 \tag{LC.12}\label{eq:lc-norm-morphism}
\]
Equations \eqref{eq:lc-jet-morphism}--\eqref{eq:lc-norm-morphism} are the
exact relation between the comparison measure and the original observation.
In particular, the original condition $p\equiv a\pmod\chi$ becomes the
condition $J_r(T_rp)=D_r^{-1}E_\chi\boldsymbol a$; it has not been discarded
or replaced by an independent arithmetic condition. The symbols
$\zeta_\lambda(r)$ describe these transported evaluation functionals only.

Let $K_N^\sigma(r)$ be the Gram matrix of all these divided-jet functionals
on $(\mathcal P_N,\|\cdot\|_\sigma)$. More explicitly, for any orthonormal
basis $e_0,\ldots,e_N$ of this space,
\[
 [K_N^\sigma(r)]_{(\lambda,j),(\mu,l)}
 =\sum_{n=0}^N
 \frac{e_n^{(j)}(\zeta_\lambda(r))}{j!}
 \overline{\frac{e_n^{(l)}(\zeta_\mu(r))}{l!}}.
 \tag{LC.13}\label{eq:lc-kernel-definition}
\]
The map $J_r$ is onto, since its restriction to $\mathcal P_{q-1}$ is
injective by the same multiplicity proof as for $E_\chi$ and the two spaces
have dimension $q$. Hence $K_N^\sigma(r)$ is positive definite.
For clarity, if $L$ is the $q$ by $N+1$ matrix of $J_r$ in an orthonormal
basis, then $K_N^\sigma(r)=LL^*$. For any required jet vector $y$, the vector
$b_0=L^*(LL^*)^{-1}y$ satisfies $Lb_0=y$. Every other solution is $b_0+h$
with $Lh=0$, and
$h^*b_0=h^*L^*(LL^*)^{-1}y=(Lh)^*(LL^*)^{-1}y=0$.
Thus the minimum squared norm is $y^*(K_N^\sigma(r))^{-1}y$.
Applying this identity to the exact original constraints yields
\[
 \boxed{
 G_N(\sigma(r\cdot),\chi)
 =r^{-1}E_\chi^*D_r^{-1}
 (K_N^\sigma(r))^{-1}D_r^{-1}E_\chi.}
 \tag{LC.14}\label{eq:lc-exact-matrix}
\]
Since $\det D_r=r^{\sum_\lambda m_\lambda(m_\lambda-1)/2}$,
\[
 \boxed{
 V_N(\sigma(r\cdot),\chi)
 =r^{-\sum_\lambda m_\lambda^2}
 \frac{|\det E_\chi|^2}{\det K_N^\sigma(r)}.}
 \tag{LC.15}\label{eq:lc-exact-determinant}
\]
This is an identity in the original remainder frame, with the original
$\chi$ on its left and the full transported jet matrix on its right.

There is an equivalent check on the remainder determinant. On
$\mathcal P_{q-1}$ the coefficient matrix $C_r$ of $T_r$ is upper triangular
with entries
\[
 (C_r)_{l,j}=\binom jl c^{j-l}(1-r^{-1})^{j-l}r^{-l}
 \quad(0\le l\le j<q),
 \qquad \det C_r=r^{-q(q-1)/2}.
 \tag{LC.16}\label{eq:lc-coefficient-map}
\]
Let $E_r$ be the divided-jet matrix at $\zeta_\lambda(r)$ on this displayed
polynomial frame. Equation \eqref{eq:lc-jet-morphism} reads
$E_rC_r=D_r^{-1}E_\chi$. Consequently
\[
 G_N(\sigma(r\cdot),\chi)
 =r^{-1}C_r^*\bigl[E_r^*(K_N^\sigma(r))^{-1}E_r\bigr]C_r,
\]
and taking determinants gives the total frame factor $r^{-q^2}$.
This explicitly agrees with \eqref{eq:lc-exact-determinant}; no modulus,
Jacobian, or derivative factor has been absorbed into a normalization.

\subsection{Passage to the asymptotic coefficient}

Write
\[
 A_\chi=\sum_{\lambda\in\Lambda}m_\lambda|\Re\lambda-c|,
 \qquad
 B_\chi=\sum_{\substack{\lambda\in\Lambda\\\Re\lambda=c}}
 m_\lambda^2.
 \tag{LC.17}\label{eq:lc-coefficients}
\]
The full confluent Gamma-kernel theorem
(Theorem~\ref{thm:fpk-full-jet-asymptotic}) proved in the preceding module,
applied to the particular fixed functionals in
\eqref{eq:lc-kernel-definition}, gives
\[
 \log\det K_N^\sigma(r)
 =rA_\chi\log N+B_\chi\log\log N+O_{r,\chi}(1).
 \tag{LC.18}\label{eq:lc-inherited-kernel-law}
\]
Indeed, $\Re\zeta_\lambda(r)-c=r(\Re\lambda-c)$, and the roots on the
line are exactly those original roots with $\Re\lambda=c$, with all their
multiplicities unchanged. This use of the Gamma theorem is for each fixed
$r>0$; no uniformity in a moving parameter is asserted or needed.
Equation \eqref{eq:lc-exact-determinant} therefore proves
\[
 \log V_N(\sigma(r\cdot),\chi)
 =-rA_\chi\log N-B_\chi\log\log N+O_{r,\chi}(1).
 \tag{LC.19}\label{eq:lc-dilated-asymptotic}
\]

We next pass the pointwise inequalities to the fixed quotient without
altering the admissible polynomial space. If $b\,w_1\le w_2\le d\,w_3$
pointwise with $b,d>0$, then for every polynomial $p$ the same inequalities
hold for the squared norms. Taking the infimum over the identical set
$\{p\in\mathcal P_N:p-a\in(\chi)\}$ gives
\[
 b\,\boldsymbol a^*G_N(w_1,\chi)\boldsymbol a
 \le\boldsymbol a^*G_N(w_2,\chi)\boldsymbol a
 \le d\,\boldsymbol a^*G_N(w_3,\chi)\boldsymbol a.
\]
Thus these are inequalities of positive definite Hermitian matrices.
If $0<H\le G$, then $H^{-1/2}GH^{-1/2}\ge I$; all its eigenvalues are
at least one, and $\det G\ge\det H$. This proves, in particular,
\[
 B_{r,M}^{-q}V_N(\sigma(r\cdot),\chi)
 \le V_N(w_M,\chi)\le V_N(\sigma,\chi).
 \tag{LC.20}\label{eq:lc-volume-sandwich}
\]
Using \eqref{eq:lc-dilated-asymptotic} at $r$ and at $1$, we obtain, for
each fixed $r>1$,
\[
 -rA_\chi\log N-B_\chi\log\log N+O_{r,M,\chi}(1)
 \le\log V_N(w_M,\chi)
 \le-A_\chi\log N-B_\chi\log\log N+O_\chi(1).
 \tag{LC.21}\label{eq:lc-log-sandwich}
\]
Dividing by $\log N$ for $N>1$ and taking lower and upper limits gives
\[
 -rA_\chi
 \le\liminf_{N\to\infty}\frac{\log V_N(w_M,\chi)}{\log N}
 \le\limsup_{N\to\infty}\frac{\log V_N(w_M,\chi)}{\log N}
 \le-A_\chi.
\]
Here $A_\chi\ge0$. The lower inequality holds separately for every fixed
$r>1$. Taking its supremum as $r$ decreases to $1$ proves the exact limit
\[
 \boxed{\displaystyle
 \lim_{N\to\infty}\frac{\log V_N(w_M,\chi)}{\log N}
 =-\sum_{\lambda\in\Lambda}m_\lambda|\Re\lambda-c|.}
 \tag{LC.22}\label{eq:lc-main-limit}
\]
This order of limits makes the possible growth of $B_{r,M}$ near $r=1$
irrelevant: the constant is fixed before $N$ tends to infinity.

If every original root lies on $\Re S=c$, then $A_\chi=0$. Already a single
fixed choice, for example $r=2$, makes the two sides of
\eqref{eq:lc-log-sandwich} equal up to bounded errors. Hence in this case
\[
 \boxed{\displaystyle
 \log V_N(w_M,\chi)
 =-\left(\sum_{\lambda\in\Lambda}m_\lambda^2\right)
 \log\log N+O_{M,\chi}(1).}
 \tag{LC.23}\label{eq:lc-critical-law}
\]
For $A_\chi>0$, equation \eqref{eq:lc-main-limit} is the conclusion proved
by this comparison; it does not assert a bounded remainder at the
$\log\log N$ scale.

\subsection{Application to the actual arithmetic weight}

Retain the actual weight $m_{h,k}$ and the original polynomial $\chi$ from
the arithmetic construction. The proved arithmetic-tail envelope for fixed
$k\ge3$ is
\[
 a_k\frac{\sigma(u)}{(1+u^2)^M}
 \le m_{h,k}(u)\le A_k\sigma(u),
 \qquad u\in\mathbb R,
 \tag{LC.24}\label{eq:lc-arithmetic-envelope}
\]
where $a_k,A_k$ are the positive constants calculated there and $M$ is its
fixed integer exponent. Applying the same quotient order argument gives
\[
 a_k^qV_N(w_M,\chi)
 \le V_N(m_{h,k},\chi)\le A_k^qV_N(\sigma,\chi).
 \tag{LC.25}\label{eq:lc-arithmetic-sandwich}
\]
The constants $k,q,a_k,A_k,M$, every root, and every root multiplicity are
fixed in this limit. Equations \eqref{eq:lc-main-limit} and
\eqref{eq:lc-dilated-asymptotic} at $r=1$ now prove
\[
 \boxed{\displaystyle
 \lim_{N\to\infty}\frac{\log V_N(m_{h,k},\chi)}{\log N}
 =-\sum_{\lambda\in\Lambda}m_\lambda|\Re\lambda-k/2|.}
 \tag{LC.26}\label{eq:lc-actual-limit}
\]
If all original roots lie on $\Re S=k/2$, the stronger conclusion is
\[
 \log V_N(m_{h,k},\chi)
 =-\left(\sum_{\lambda\in\Lambda}m_\lambda^2\right)
 \log\log N+O_{h,k,\chi}(1).
 \tag{LC.27}\label{eq:lc-actual-critical-law}
\]
These statements concern the actual fixed quotient. They provide no
uniform estimate when the arithmetic polynomial or $k$ changes with $N$,
and they do not determine a signed difference of growing-window quotient
determinants. All such parameters remain in their original domains.

Finally retain the original tensor-Gamma density $m_0$ of (AT4),
with its original mass $(2\pi)^{k/2}$. Its proved complete asymptotic
(FPK.29) has the same leading coefficient $-A_\chi$. Hence the
actual endpoint quotient ratio satisfies
\[
 \lim_{N\to\infty}\frac{1}{\log N}
       \log\frac{V_N(m_{h,k},\chi)}{V_N(m_0,\chi)}=0.
 \tag{LC.28}\label{eq:lc-original-reference-ratio}
\]
If all original roots lie on $\Re S=k/2$, (LC.27) and (FPK.29)
give the stronger bounded logarithmic ratio, equivalently positive
upper and lower bounds for this ratio for all sufficiently large $N$.
Both conclusions are for the original fixed $h,k,\chi$. No
uniformity needed for a varying-packet four-window expression is
asserted by (LC.28).
